% ===========================================================================
% KALİBRASYON BÖLÜMÜ REVİZYON NOTLARI (§3.3, §3.4, §4.1, Appendix B/C/F)
% Bu dosya Overleaf'e doğrudan girmez; mevcut metinde yapılacak değişiklikleri
% "ESKİ → YENİ" blokları halinde toplar. Yeni paragraflar kopyala-yapıştır
% için hazırdır.
% ===========================================================================

% ---------------------------------------------------------------------------
% 1) CONFIG YENİDEN NUMARALANDIRMA
% C2'nin pilotta çıkarılmış olması nedeniyle metin C0, C1, C3...C11 diye
% ilerliyordu; tüm configler ardışık yeniden adlandırılır. Eşleme:
%
%   eski C0  -> C0   (demografik baseline)
%   eski C1  -> C1   (tam profil)
%   eski C3  -> C2   |
%   eski C4  -> C3   |
%   eski C5  -> C4   |  leave-one-group-out varyantları: C2-C7
%   eski C6  -> C5   |
%   eski C7  -> C6   |  <-- SEÇİLEN KONFİGÜRASYON (her iki nihai protokol)
%   eski C8  -> C7   |
%   eski C9  -> C8   |
%   eski C10 -> C9   |  indirgeme varyantları: C8-C10
%   eski C11 -> C10  |
%
% Yapılacaklar:
%  - §3.3'teki "labeled C0, C1, and C3 through C11" cümlesi ->
%    "labeled C0 through C10" olur.
%  - Dipnot 2 ("C2 was reserved for an earlier pilot configuration...")
%    TAMAMEN SİLİNİR.
%  - Appendix B Table 4, Appendix C Tablo 5-12 ve §4.1 metnindeki tüm config
%    etiketleri eşlemeye göre güncellenir (metin içinde: C7->C6, C11->C10,
%    C3->C2, C4->C3, C8->C7, C9->C8, C10->C9).
% ---------------------------------------------------------------------------

% ---------------------------------------------------------------------------
% 2) NİHAİ PROTOKOL DEĞİŞİKLİĞİ
% Alınan karar: her iki outcome için persona konfigürasyonu C6 (eski C7).
%   - Binary (pacdemons):  C6 + direct answering + T=0   + GPT-4o-mini
%   - Likert (womenwork):  C6 + verbalized sampling + T=0.8 + Gemini 2.5
%                          Flash Lite
% Mevcut taslak pacdemons için T=0.8 ve womenwork için C11+GPT-4o-mini
% seçiyordu; §4.1'in protokol seçim paragrafları ve Tablo 1-2, Appendix F
% buna göre revize edilmelidir. Bilinen hücre değerleri (mevcut Appendix C
% tablolarından, eski numaralarla):
%   pacdemons  C7 + D + T=0   : JSD 0.001, MCC 0.410, kappa 0.405, Rec+ 0.512
%   womenwork  C7 + VS + T=0.8 + Gemini : Wass 0.218, JSD 0.030,
%                                         kappa 0.151, kappa_w 0.401
% DİKKAT: Appendix F'teki "Gemini 2.5 Flash Lite at C7 ... 0.47 and 1.33
% across two additional seeds" cümlesi yeni seçimle çelişir; bu paragraf ya
% güncel çok-tohumlu koşu sonuçlarıyla değiştirilmeli ya da çıkarılmalıdır.
% ---------------------------------------------------------------------------

% --- YENİ protokol seçim paragrafı (§4.1 sonundaki "We therefore fix..."
%     paragrafının yerine):

We therefore fix a single persona configuration, C6, for both outcomes, and
let the inference protocol adapt to the response format. For the rare binary
participation outcome, the final protocol is C6 with direct answering at
$T=0$, which preserves respondent-level discrimination on the rare
affirmative class. For the ordered Likert outcome, the final protocol is C6
with verbalized sampling at $T=0.8$ on Gemini~2.5 Flash Lite, which offers
the strongest joint performance when aggregate ordinal fit and
respondent-level ordinal agreement are considered together: its quadratic
weighted kappa is roughly five times that of the best distribution-only
alternative at a Wasserstein distance that remains close to the minimum
observed in the search. Fixing one persona configuration across outcomes
also simplifies the experimental stage, where both outcome formats are
administered to the same synthetic population.

% --- YENİ recall gerekçelendirme cümleleri (§4.1'de pacdemons protokolü
%     anlatılırken; "recovering 53.6% of true affirmative cases" tarzı
%     cümlelerin yerine — T=0 kararıyla tutarlı %51,2 kullanılmıştır):

Recall on the rare affirmative class deserves emphasis rather than apology:
the calibrated protocol identifies roughly half of the respondents who
actually attended a protest in the previous twelve months
(protest-positive recall $= 0.51$), while holding the simulated participation
rate within two percentage points of the observed five percent. In other
words, the protocol gains recall without paying for it in precision, where
the characteristic failure mode of persona-conditioned models on rare
political behavior is to collapse onto the majority class and find almost no
true participants at all, as the demographic-only baseline does
(recall $= 0.01$).

% --- YENİ womenwork gerekçelendirme cümleleri (§3.2'de womenwork'ün
%     tanıtıldığı paragrafa eklenecek):

\emph{Womenwork} is selected as the attitudinal calibration target for three
reasons. Its five-point Likert structure matches the response format of the
evaluative outcomes in the downstream experiment, so the protocol selected
on it transfers directly. Its subject matter, gendered divisions of family
labor, indexes orientation toward progressive social change while remaining
distinct from the experimental domain of Kurdish-language rights, which
guards against calibrating on the very attitudes the experiment
manipulates. And its observed distribution is genuinely graded rather than
consensual, so it provides a demanding test of whether a protocol can
recover an ordered distribution rather than a single modal answer.

% --- YENİ model gerekçelendirme cümlesi (§3.4 "Cross-model robustness"
%     paragrafına, model listesinin ardına eklenecek):

The four backbones are chosen because they are among the most widely used
commercial and open-weight model families in applied synthetic-survey work,
spanning three providers and both proprietary and open-weight release
models; this anchors the robustness check in the models practitioners
actually deploy rather than in an exhaustive sweep of the model space.

% ---------------------------------------------------------------------------
% 3) METRİK YOĞUNLUĞUNU AZALTMA (uygulama talimatı)
%  - §3.2'deki metrik tanıtım paragrafı korunur ama §4.1 ana metninde her
%    iddia için EN FAZLA BİR metrik verilir; ikinci metrikler Appendix C
%    tablolarına bırakılır.
%  - §4.1'deki dört bold bulgu başlığı AYNEN KORUNUR (yapı değişmiyor);
%    başlık altı paragraflarda sayı yoğunluğu azaltılır: örn. "minority-class
%    recall ranges from 0.01 ... to 0.51" gibi tek karşıtlıklar kalır,
%    ara değer sıralamaları (ör. üç sıcaklık değerinin tamamı) tablolara
%    devredilir.
%  - Kappa/QWK/MCC kısaltmaları ilk geçtiği yerde bir kez açılır, sonra
%    yalnızca kısaltma kullanılır.
% ---------------------------------------------------------------------------

% ---------------------------------------------------------------------------
% 4) FUATKINA FEEDBACK'İNDEN GELEN EK DEĞİŞİKLİKLER (PDF yorumları)
%  - "minority-class recall" -> metin ve tablolarda "protest-positive
%    recall"; Macro-F1'in de protest-pozitif sınıfa göre tanımlandığı tablo
%    başlığında not edilir. (s.15)
%  - §3: "Empirical Strategy" üst başlık olur; Data, Persona Construction,
%    Calibration Procedure onun alt başlıkları; girişe iki aşamayı ve bölüm
%    planını anlatan kısa yol haritası paragrafı eklenir. (s.8)
%  - Distributional-alignment gerekçe pasajı Fuat'ın referanslı metniyle
%    değiştirilecek — metni Fuat'tan al. (s.9)
%  - Appendix D: persona tablosunun altına kalibrasyonda kullanılan tam
%    prompt (system + user) eklenir; deney promptları zaten
%    appendix_experiment_materials.tex'te. (s.10)
%  - "Configuration" terimi daha betimleyici bir adla değiştirilir (öneri:
%    "persona content configuration" ilk kullanımda tanımlanıp kısaltılır);
%    Figure 1 bölümün sonuna taşınır; funnel, tam grid'in maliyeti ve toplam
%    kombinasyon sayısı verilerek gerekçelendirilir. (s.10)
%  - Verbalized sampling'e referans eklenir. (s.11)
%  - T ile sampling stratejisinin neden birlikte tarandığına dair bir cümle
%    eklenir. (s.11)
%  - Metin içi "Appendix X" referanslarının çoğu footnote'a taşınır. (s.11)
%  - Tablo 1'den CV sütunu çıkarılır; T sütunundan sonra dikey çizgi. (s.14)
%  - Tablo 14'te "Error" sütun adı "Difference" olur. (s.40)
%  - §4.1'in dili sadeleştirilir (teknik terim yoğunluğu azaltılır). (s.16)
% ---------------------------------------------------------------------------

% ---------------------------------------------------------------------------
% 5) APPENDIX YENİDEN DÜZENİ + NATURAL REWRITE ENTEGRASYONU
%  Yeni appendix akışı:
%    A  Persona Variables                  (değişmez)
%    B  Persona Configurations             (yeni C-numaralamasıyla)
%    C  Calibration Funnel: Cell-Level Results
%       C.1-C.3                            (değişmez, yeni numaralama)
%       C.4 Prompt-framing stress tests    (+ ideoloji arka plan paragrafı
%            TR+EN: appendix_C4_background_paragraph.tex, \label{app:background})
%       C.5 Persona format ablation        (appendix_natural_ablation.tex)
%    D  Anonymized Persona Example         (+ tam kalibrasyon promptu:
%            system + user; Fuat feedback s.10)
%    E  Verbalized Sampling Prompt         (değişmez; C.5 ve H buna referans
%            verir, metin tekrarı yapılmaz)
%    F  Stochastic Stability               (Gemini cümlesi revize — madde 2)
%    G  Subgroup Validation                (değişmez)
%    H  Experimental Materials             (appendix_experiment_materials.tex)
%
%  §3.4 "Prompt-framing stress tests" paragrafı dört varyantlı revize haliyle
%  değiştirilir (natural rewrite dördüncü varyant; footnote'lar
%  \ref{app:background} ve \ref{app:natural}'a işaret eder).
%
%  §4.1'deki stres testi bulgu paragrafının sonuna eklenecek cümle:
%    "The natural-language rewrite is the exception that clarifies the rule:
%     presentation choices that leave the structured variable map intact
%     barely move the results, whereas dissolving that map into prose halves
%     respondent-level agreement (Appendix~\ref{app:natural})."
% ---------------------------------------------------------------------------
